\documentclass[10pt]{article}
\usepackage[margin=20mm]{geometry}
\usepackage{amsmath,amssymb,mathtools,fontspec,unicode-math,xurl,hyperref}
\IfFontExistsTF{DejaVu Serif}{\setmainfont{DejaVu Serif}}{\setmainfont{Latin Modern Roman}}
\setmathfont{latinmodern-math.otf}
\setlength{\parindent}{0pt}\setlength{\parskip}{6pt}\sloppy
\title{Technical Verification of Mack Modes in Supersonic Boundary Layer}
\author{AI-assisted technical review}
\date{}
\begin{document}
\maketitle
Abstract: Two core derivations in the acquired preprint require repair. The stated Airy functions solve an operator whose potential differs from the printed potential by a nonzero explicit residual. The resulting Green formula therefore does not invert the printed operator, and the subsequent exact gluing identities do not follow. Separately, the positive critical-layer imaginary-part bound uses a curvature nondegeneracy that is absent from the structural assumptions. An analytic monotone profile realizes the omitted degenerate case. These are Major proof defects in the reviewed version; neither finding establishes that the final instability theorems are false.\par
\section*{Identity, scope and operative hypotheses}
The reviewed paper is N. Masmoudi, Y. Wang, D. Wu and Z. Zhang, Mack modes in supersonic boundary layer, DOI 10.1002/cpa.70022 (https://doi.org/10.1002/cpa.70022). The acquired source is arXiv:2405.04853v1, 8 May 2024, 91 pages [paper]. Physical and printed page numbers agree. Journal metadata records online publication on 20 November 2025 and volume 79 (2026), 897–1011. This report concerns the exact 91-page preprint, without asserting equivalence to the journal article.\par
The root reviewer read source pp. 1–32, 36–37, 39–52, 57–62 and 79–91. The review directly checks the pressure reduction, Airy operator, critical-layer nonvanishing mechanism and their dependence in the final conclusions. Pages 33–35, 38, 53–56 and 63–78 were not fully read. The complete mechanical structure index is retained separately and is not represented as scientific reading. The findings below are proved within this declared scope.\par
Write \(T=T_0\), \(c=c_r+ic_i\), and \(M=M_a(U_B-c)/\sqrt T\), suppressing the source’s bar on \(M\). The structure condition (1.4), p. 3, requires a smooth strictly increasing profile from \(0\) to \(1\), with exponential decay of \(1-U_B\) and its first three derivatives. The temperature is \[T=1+\frac{\gamma-1}{2}M_a^2(1-U_B^2),\qquad \gamma=1.4.\] The construction uses \(M_a>1\), a simple turning point \(Y_0\), a separated critical point \(U_B(Y_c)=c_r\), and large \(\alpha\) with \(c_i>0\) smaller than every inverse power in the regime \(\mathbb H_0\), p. 13. Cutoff distances are independent of \(\alpha\); the inner interval lies where \(F_r=1-M_r^2>0\). Theorem 1.1, p. 4, also requires \(c_r\in\mathcal H(U_B,M_a)\cap(1-M_a^{-1},1)\), where \(\mathcal H\) is the nonzero-\(J\) set. Analyticity is an additional hypothesis for Theorem 1.4 at sufficiently large Mach number. The formal Blasius statement, Theorem 1.5, p. 5, is for \(M_a\geq3\). The discussion of \(2.2\) is not treated as a proved extension.\par
\section*{The Airy operator and the gluing identities}
Equation reference: eq:operator, eq:printed, eq:airy, eq:langer. Equations (3.1)–(3.4), p. 12, and (3.7)–(3.12), p. 13, define \[\begin{aligned}
 L_p P&=P''-2(M'/M)P'-\alpha^2\widetilde F P-Q_pP,
 \\
 Q_p&=\frac{M''}{M}-2\left(\frac{M'}M\right)^2
 -\frac34\left(\frac{\eta''}{\eta'}\right)^2
 -\frac{\eta'''}{2\eta'},\\
 E&=M(\eta')^{-1/2},\quad
 A=E\operatorname{Ai}(\kappa\eta),\quad
 B=E\operatorname{Bi}(\kappa\eta),\\
 \kappa^3&=\alpha^2F_r'(Y_0),\qquad
 F_r'(Y_0)\eta(\eta')^2=\widetilde F,\qquad \eta(Y_0)=0.
 \end{aligned}\] The source asserts \(L_p A=L_p B=0\).\par
Problem 1 (P01: Major; invalid exact outer inverse).  The functions in [eq:airy] (\#eq:airy) do not in general solve the printed homogeneous equation. They have the exact residual \[L_p A=RA,\qquad L_p B=RB,\qquad
 R=\frac32\left(\frac{\eta''}{\eta'}\right)^2.\] Consequently the printed Green formula and the exact telescoping cancellation in the gluing argument require correction.\par
Equation reference: eq:corrected. Proof. For either Airy function \(a\), differentiate \(Ea(\kappa\eta)\). The coefficient of \(\kappa a'\) is \(2E'\eta'+E\eta''-2(M'/M)E\eta'=0\). The Airy equation and [eq:langer] (\#eq:langer) cancel the \(\alpha^2\widetilde F\) term. The remaining coefficient, divided by \(E\), is \[Q_c=\frac{E''-2(M'/M)E'}{E}
 =\frac{M''}{M}-2\left(\frac{M'}M\right)^2
 +\frac34\left(\frac{\eta''}{\eta'}\right)^2
 -\frac{\eta'''}{2\eta'}.
 \] Thus \(Q_c-Q_p=R\), proving the claim without an asymptotic approximation. ◻\par
\subsection*{An admissible nonzero residual}
Take \(M_a=3\), \(c_r=9/10\), \(U_B=1-e^{-Y}\) and, for example, \(c_i=e^{-\alpha}>0\). This analytic monotone profile satisfies (1.4). With \(u=U_B\), \[T=(14-9u^2)/5,\quad
 F_r=1-\frac{9(u-9/10)^2}{T},\quad
 u_0=\frac34-\frac{\sqrt{4755}}{180}.\] Here \(F_r(0)=-449/280<0\), \(F_r(\infty)=91/100>0\); \(Y_0\simeq0.45714\), \(Y_c=\log10\), and the fixed gap is about \(1.84545\). Near \(Y_0\), the cutoff gives \(\widetilde F=F_r\). Taylor expansion of [eq:langer] (\#eq:langer) gives \(\eta'(Y_0)=1\) and \(\eta''(Y_0)=F_r''(Y_0)/(5F_r'(Y_0))\). Exact algebra yields \[\frac{F_r''(Y_0)}{F_r'(Y_0)}
 =-\frac{1253158}{450241}
  +\frac{1624005\sqrt{4755}}{142726397},
 \qquad
 R(Y_0)=\frac3{50}\left(\frac{F_r''(Y_0)}{F_r'(Y_0)}\right)^2
 \simeq0.2396849464>0.\] The radical expression is nonzero exactly; the decimal is only illustrative. Since \(M(Y_0)\ne0\) and \(\operatorname{Ai}(0)>0\), the residual \(RA\) is nonzero. No assertion \(J(9/10;3)\ne0\) is needed for the Section 3 identity, and none is made.\par
\subsection*{Dependency, alternate reading and repair}
The Wronskian stated on p. 13 remains correct. It gives a Green operator \(G\) for the potential \(Q_c\), so the same formula satisfies \[L_pG[f]=f+RG[f],\] rather than \(f\). The error is then visible at the initial gluing step, p. 44: the asserted decomposition \(L[G_1(A)]=H_1[A]+H_2[A]\), with (5.7), p. 46, omits the extra \(R\chi_1 A\). The subsequent cancellations, pp. 45–46, and the proof of Proposition 5.1, p. 52, use those identities. This affects the derivation of the boundary expansion in Proposition 6.1 and hence the dispersion relation and Theorem 7.1, pp. 79–80.\par
Replacing the minus sign in [eq:printed] (\#eq:printed) by a plus sign restores the Airy identity. Away from the critical layer the added term has the same admissible qualitative decay as the existing potential; this is evidence that a repair may be possible. However, the printed minus sign is also present in \(\widetilde Q_1\), (6.3), p. 57, and in the calculation on p. 83. The coefficient \(J\) in (6.2), the set \(\mathcal H\) in (6.4), and the lower bound deduced from (6.15), p. 60, must therefore be reconsidered. Treating all subsequent formulas as if they silently used \(Q_c\) is not a verified alternate reading. A repaired construction must propagate the change and establish its nonvanishing conditions. This Major classification concerns an invalid core derivation, not a demonstrated absence of Mack modes or a false final theorem.\par
\section*{A missing critical-curvature hypothesis}
Lemma 6.2, p. 58, assumes \((\alpha,c)\in\mathbb H_0\) and cutoffs as in Lemma A.2. For sufficiently small positive \(c_i\), it asserts a nonzero imaginary contribution with \[|\operatorname{Im}\mathcal K_B[\varphi_0]|
 \asymp \alpha^{5/6}e^{-\alpha w_0(\mathring Y)},
 \qquad \mathring Y\in[Y_c,\widetilde Y].\] Its hypotheses do not require \(M_r''(Y_c)\ne0\). Lemma 4.7, pp. 41–43, explicitly retains the factor \[\operatorname{Im}\int\frac{F}{\phi^2}
 =-\frac{\pi M_r''(Y_c)}{M_r'(Y_c)^3}
   +O(\alpha^2c_i|\log c_i|).\] In the proof of Lemma 6.2, p. 90, the leading term is proportional to \(M_r''(Y_c)\), before a positive lower bound is inferred with this factor absent.\par
Problem 2 (P02: Major; unjustified nonvanishing).  Strict monotonicity does not imply the critical curvature required by this lower-bound argument. A profile satisfying the structural assumptions can have \(M_r''(Y_c)=0\). The inference to a uniform positive imaginary-part lower bound then fails. The degenerate inner equation also has a real regular limit, which excludes the stated lower bound for arbitrarily small \(c_i>0\) with fixed admissible cutoffs.\par
Equation reference: eq:profile. Profile and limiting argument. Keep \(\gamma=7/5\), \(M_a=3\), \(c_r=9/10\), and set \(a=5090/671\). Define \(U_B\) by \[Y=\int_0^{U_B(Y)}\frac{e^{-as}}{1-s}\,ds,
 \qquad U_B'=(1-U_B)e^{aU_B}.
 \] The positive integrand gives an analytic increasing inverse. The integral diverges at \(1\), so the profile is defined on the whole half-line and increases from \(0\) to \(1\). As \(Y\to\infty\), \(1-U_B\) and its first three derivatives decay exponentially, with limiting rate \(e^a\); choosing a smaller positive exponent gives precisely the uniform bounds required by (1.4). The turning velocity is the same \(u_0<c_r\) as above. Its distance from the critical point is positive and independent of \(\alpha\). A fixed rescaling of \(Y\), if desired, normalizes this distance without changing the curvature cancellation.\par
Put \(m=M_r=M_a(U_B-c_r)/\sqrt T\). At \(Y_c\), \(T_c=671/500\), and direct differentiation gives \[\begin{aligned}
 m''(Y_c)
 &=\frac{M_a}{\sqrt{T_c}}\left[
 U_B''(Y_c)+
 \frac{(\gamma-1)M_a^2c_r}{T_c}(U_B'(Y_c))^2\right],\\
 \frac{U_B''(Y_c)}{(U_B'(Y_c))^2}
 &=-\frac1{1-c_r}+a=-\frac{1620}{671},
 \qquad
 \frac{(\gamma-1)M_a^2c_r}{T_c}=\frac{1620}{671}.\end{aligned}\] Thus \(m''(Y_c)=0\) exactly, while \(m'(Y_c)>0\).\par
Equation reference: eq:l1, eq:inner. Here is a direct check of the limiting mechanism, independent of the asymptotic inference on p. 90. Write \(M_\epsilon=m-i\epsilon b\), \(b=M_a/\sqrt T\), \(\epsilon=c_i>0\), and \(F_\epsilon=1-M_\epsilon^2\). Because \(m\) has a simple zero and \(m''(Y_c)=0\), \(a_0=m''/m\) extends analytically through \(Y_c\). On a fixed sufficiently small inner interval \(I\) about \(Y_c\), \(F_0=1-m^2\) is bounded below by a positive constant. Moreover, \[\frac{M_\epsilon''}{M_\epsilon}
 =a_0+\frac{i\epsilon(a_0b-b'')}{m-i\epsilon b}
 \longrightarrow a_0\quad\hbox{in }L^1(I),
 \] with error \(O(\epsilon|\log\epsilon|)\). Dividing the inner equation (4.1), p. 28, by \(M_\epsilon\) and multiplying by \(F_\epsilon\) gives \[v''-\frac{F_\epsilon'}{F_\epsilon}v'
 -\left(\alpha^2F_\epsilon+
 \frac{M_\epsilon''}{M_\epsilon}
 +\frac{2(M_\epsilon')^2}{F_\epsilon}\right)v
 =\frac{F_\epsilon f_\epsilon}{M_\epsilon},
 \qquad v|_{\partial I}=0.
 \] For fixed sufficiently large \(\alpha\), the real limit operator is invertible: after division by \(F_0\), integration by parts leaves the positive derivative energy and an \(\alpha^2\) term dominating all bounded potential terms. Thus the homogeneous Dirichlet problem has only the zero solution.\par
The forcing of Lemma 6.2 is \(f_\epsilon=\partial_Y(-H_1[A]/F_\epsilon)\); see (5.7), p. 46, and (5.13)–(5.14), p. 48. It vanishes in a neighborhood of \(Y_c\): there \(\chi_1=0\), \(\chi=1\), and \(H_1[A]/F_\epsilon\) is a constant integral. Away from \(Y_c\) it converges to real smooth forcing. Consequently the right side of [eq:inner] (\#eq:inner) converges to a real function. Equation [eq:l1] (\#eq:l1) and the remaining uniform coefficient convergences imply convergence of the first-order fundamental matrices by their integral equations and Gronwall’s inequality. The nonzero limiting Dirichlet determinant then gives convergence of \(v_\epsilon\) to the real Dirichlet solution.\par
The functional \(\mathcal K_B\), (5.8), p. 46, has kernel supported where the cutoff derivatives lie, away from \(Y_c\). Its coefficients converge there. Hence \(\operatorname{Im}\mathcal K_B[v_\epsilon]\to0\) as \(\epsilon\downarrow0\) for each fixed sufficiently large \(\alpha\). Cutoffs and asymmetric endpoints can be placed inside \(I\) with the fixed separations and \(F_r\geq9/10\) used in Lemma A.2, p. 86, by choosing the separation constants sufficiently small. Their geometry does not force \(m''(Y_c)\ne0\). A lower bound \(C^{-1}\alpha^{5/6}
e^{-\alpha w_0(\widetilde Y)}>0\), uniform in arbitrarily small \(\epsilon\), is therefore impossible in this case. For the joint large-\(\alpha\) regime one can choose a diagonal \(\epsilon(\alpha)>0\) below every inverse power and below the required exponential threshold, with the imaginary part smaller than this proposed bound by an arbitrarily small factor. ◻\par
\subsection*{Impact and repair boundaries}
The curvature factor is already visible in the source’s own leading-term calculation, so the issue is not merely a different convention for the sign of a complex quantity. A lower bound can depend on a fixed nonzero \(|m''(Y_c)|\); it cannot absorb a zero factor into a positive implicit constant. Lemma 6.2 feeds \(\operatorname{Im}\mathfrak D_0\) in Lemma 7.3, p. 80, and the positive growth rate at the end of Theorem 7.1, p. 83. Thus the missing hypothesis is material.\par
One possible repair is an explicit nondegeneracy condition, carried through the parameter set and main theorem. Another is a separate analysis of the degenerate critical layer, including the true size and sign of its contribution. Correcting P01 does not supply this condition. The profile [eq:profile] (\#eq:profile) is not asserted to satisfy \(J(c_r;M_a)\ne0\). Accordingly this report proves a defect in the general lower-bound argument, not a counterexample to the final conditional theorem or to the Blasius-specific assertion. The mechanism may remain valid on a suitable nondegenerate subset.\par
\section*{Independent checks and version limits}
The pressure equation itself can be obtained from (1.5), p. 3. With \(q=U_B-c\), consistent reconstruction formulas are \[\begin{aligned}
 V&=\frac{iTP'}{\gamma M_a^2\alpha q},&
 U&=-\frac{TP}{\gamma M_a^2q}
     -\frac{TP'U_B'}{\gamma M_a^2\alpha^2q^2},\\
 \rho&=\frac{P}{\gamma M_a^2q^2}
 +\frac{2U_B'P'/q^3-P''/q^2}{\gamma M_a^2\alpha^2},&
 \mathcal T&=PT-\rho T^2.\end{aligned}\] Substituting \(P''=(2U_B'/q-T'/T)P'
+\alpha^2(1-M_a^2q^2/T)P\) makes all four linearized residuals vanish. These identities, the Airy residual, and both explicit profiles were checked by independently written symbolic calculations, executed with SymPy 1.13.3; exact results and execution receipts are retained. No author code was executed and no numerical spectral computation is claimed.\par
The printed density reconstruction in (2.4), p. 6, needs algebraic correction. Also the complex increment in (2.7) must have the coefficient \(2ic_i(U_B-c_r)\); the calculation on p. 59 already uses this factor two. These are not promoted to additional substantive Problems, since the direct reduction and the later use supply the necessary correction. The Wronskian and the Airy scale \(\kappa\asymp\alpha^{2/3}\) are consistent. For a normal mode \(e^{i\alpha X-i\alpha ct}\), the time growth is \(e^{\alpha c_i t}\). The announced \(c_i\asymp\alpha^{-3}e^{-C\alpha}\) gives positive but vanishing growth rate \(\alpha c_i\), not a uniform high-frequency rate or a nonlinear instability conclusion.\par
The current arXiv metadata retrieved for this work lists v1 as the available version. The publisher request returned HTTP 403; journal full text was not compared. The retrieved Crossref record did not provide a correction relation. These observations do not certify that no correction exists. External proofs cited by the paper were not all acquired and independently reproved; they are not needed for the two explicit defects above. No missing executable material is classified as a mathematical error.\par
\section*{Disposition}
P01 and P02 require substantive proof repair in the reviewed preprint. The report distinguishes the false exact Airy identity and the invalid nonvanishing argument from the truth of the final instability conclusions. This initial technical review is complete in the declared scope; scientific decision remains pending and independent acceptance has not been performed.\par
9 N. Masmoudi, Y. Wang, D. Wu and Z. Zhang, Mack modes in supersonic boundary layer, arXiv:2405.04853v1 (8 May 2024), https://arxiv.org/abs/2405.04853v1. Journal identifier: https://doi.org/10.1002/cpa.70022.\par
\end{document}
